\documentclass[11pt]{article}

\usepackage[margin=1in]{geometry}
\usepackage[utf8]{inputenc}
\usepackage[T1]{fontenc}
\usepackage{lmodern}
\usepackage{amsmath,amssymb,amsthm}
\usepackage{booktabs}
\usepackage{tabularx}
\usepackage{microtype}
\usepackage[numbers,sort&compress]{natbib}
\usepackage{xcolor}
\usepackage[colorlinks=true,linkcolor=blue!50!black,citecolor=blue!50!black,urlcolor=blue!50!black]{hyperref}

\newcommand{\draftnote}[1]{\textcolor{red!70!black}{[\textbf{Draft:} #1]}}

\title{How Large a Machine, Told How Much?\\The Learned Machines' Sizes Against the Literature\\[0.5em]
\large Research note, 27 October 2026: not for circulation}

\author{
  David Barkhuizen\\
  Independent researcher
  \and
  Claude (Opus 5.5)\\
  Anthropic
}

\date{}

\begin{document}

\maketitle

\draftnote{Written at the first author's request on 27 October 2026. Our numbers are the registered
findings in \texttt{indrajala-ml}'s \texttt{docs/\allowbreak workplans/\allowbreak paper-workplan.md}
(at b97267b: T2, T5, T10 and the approximation claim) and paper~1's draft at 44f87d9. \textbf{T11 and
T10r are not read}: T11's long tests and T10r's batch are running, and none of their data is used
here. The literature is the repository's \texttt{library.bib} with the papers re-read for their
machines' sizes, plus a search on 27 October that added
\citet{giles-1995-learning-class-large-finite-state-machines} and
\citet{aichernig-2024-learning-minimal-automata-rnn}.}

\section{The claim, and why it needed correcting}

We had said, in conversation and in the paper workplan's statement of the result, that our networks
learn larger machines, proved correct at every length, than any reported: finite-state transducers of
100--169 states, where ``the extracted-machine literature stops at a few states''. Read as stated, the
claim is \textbf{false}. Two works report larger machines learned by gradient descent and checked
exactly:
\begin{itemize}
\item \citet{brandoit-2026-length-independent-state-tracking} learn the word problems of $S_4$ (24
  states), $A_5$ (60) and the Mathieu group $M_{11}$ (7{,}920), on all 5 of 5 seeds each. Every
  trained model's transition tables are enumerated from the initial joint index over every symbol, so
  the check is exhaustive and proves the model correct at every length. Paper~1's draft cites this
  work, but not its machines' sizes.
\item \citet{giles-1995-learning-class-large-finite-state-machines} learn finite memory machines of
  256 and 512 states with a recurrent network that has no hidden state. At about 250 training strings
  the 512-state machine is learned with zero error on the complete set of 2{,}046 strings that
  identifies it, and the extracted machine is minimised.
\end{itemize}
What separates both from our largest results is \emph{how much of the state the training signal
tells the network}, not the size. The corrected claim therefore needs that variable.

\section{Terms}

A network is \emph{told the state} when its training signal determines the target's state at every
step. That happens in three ways:
\begin{enumerate}
\item it is trained on the state itself, as with \citeauthor{brandoit-2026-length-independent-state-tracking}'s
  labels (``the state supervised at every position'');
\item it is trained on outputs that jointly determine the state, as with division given its
  cover, whose sum and difference determine $(A \bmod p, B \bmod p)$ for odd $p$; or
\item the state is a function of a bounded window of past inputs and outputs that the network
  observes, as in a finite memory machine whose past outputs are fed back by teacher forcing
  \citep{giles-1995-learning-class-large-finite-state-machines}.
\end{enumerate}
Otherwise the state is \emph{hidden}: the network sees only outputs that leave most of the state
undetermined, so the outputs alone do not fix it. A machine is \emph{proved} when a complete
enumeration shows the network equal to the target on every input of every length. It is
\emph{tested} when it is checked on samples, test sets or bounded lengths, which can leave a wrong
network undetected (paper~1's K2: 2 of 105 addition networks passed every sampled test and are wrong).

Our runs give a state bound but no structure: one cell for every target, its codes at least three
times the target's states. On division and T11, the cell's two factors match the two coordinates
given.

\section{The literature, by what it is told and how it is checked}

\begin{table}[h]
\centering
\small
\begin{tabularx}{\textwidth}{Xlrl}
\toprule
work & signal & largest machine & check \\
\midrule
\citet{zeng-1993-self-clustering-recurrent-networks} & accept labels (hidden) & 4 & extracted, minimised, equal \\
\citet{tino-1995-learning-extracting-initial-mealy-automata} & Mealy outputs (hidden) & 5 & extracted \\
\citet{wang-2019-state-regularized-rnns} & accept labels (hidden) & Tomita, $\le 5$ & minimal DFA recovered \\
\citet{koul-2019-finite-state-recurrent-policies} & policy (hidden) & 3 memory states & extracted after training \\
\citet{aichernig-2024-learning-minimal-automata-rnn} & Mealy outputs (hidden) & 5 (minimal) & isomorphism with the target \\
\citet{umili-2024-deepdfa-automata-learning-probabilistic-relaxations} & accept labels (hidden) & 30 (random DFAs) & \textbf{tested} (to length 150) \\
\citet{weiss-2018-extracting-automata-queries-counterexamples} & (extraction of trained RNNs) & targets of 10 & extracted machines approximate \\
\citet{papazov-2025-exact-learning-arithmetic-differentiable-agents} & expert trajectories (told) & 108 & \textbf{tested} \\
\citet{giles-1995-learning-class-large-finite-state-machines} & outputs, fed back (told, by window) & 512 & complete identifying set \\
\citet{brandoit-2026-length-independent-state-tracking} & state labels (told) & 7{,}920 & \textbf{proved} (exhaustive tables) \\
\bottomrule
\end{tabularx}
\caption{Gradient-trained networks that learn a finite-state machine, with their largest target in
states.}
\label{tab:literature}
\end{table}

Two of these call for care. \citeauthor{giles-1995-learning-class-large-finite-state-machines}'s
network taps are set to the machine's input and output orders (5 and 4), so the architecture is given
the machine's memory structure. Their zero error is an average over runs at the larger training sets,
not a per-network certificate, and the 256-state machine was ``not usually'' learned perfectly.
\citeauthor{brandoit-2026-length-independent-state-tracking}'s heads are sized to the group's action
(on $M_{11}$, 4 heads of 11 attractor indices, each tracking where one point is sent), so the
architecture is given a factorisation. Both are larger than ours all the same, and both are learned
by gradient descent and checked exactly or completely.

\section{Ours, by the same measure}

\begin{table}[h]
\centering
\small
\begin{tabularx}{\textwidth}{>{\raggedright}X>{\raggedright}Xlr>{\raggedright\arraybackslash}X}
\toprule
target & extra outputs & state & $q$ & certified (registered) \\
\midrule
addition (carry automaton) & none & hidden & 3 & 52 of 105 networks \\
division, $p = 5$ & none & hidden & 25 & 1 of 20 \\
division, $p = 11$ & none & hidden & 121 & 0 of 20 (T2) \\
$S_3$ word problem & none & hidden & 6 & 1 of 20 (T5) \\
$S_4$ word problem & none & hidden & 24 & 4 of 20 (T5) \\
move-to-front, 4 items & none & hidden & 24 & 4 of 20 (T5) \\
random machine & none & hidden & 12 & 3 of 20 (T5) \\
random machine & none & hidden & 16 & 0 of 20 (T5) \\
random, primitive & none & hidden & 16 & 2 of 15 machines, 2 of 30 runs (T10, P0) \\
random, planted 4 blocks & none & hidden & 16 & 3 of 15 machines (T10, P1) \\
random, planted 4 blocks & the planted quotient & partly told & 16 & 8 of 15 machines (T10, P2) \\
division, $p = 5$ & sum & nearly told & 25 & 10 of 20 \\
division, $p = 5, 7, 11, 13$ & sum, difference & told & 25--169 & 15, 17, 16, 15 of 20 (T2) \\
$S_3$, $S_4$, move-to-front & a pair that determines the state & told & 6, 24 & T11, not read \\
\bottomrule
\end{tabularx}
\caption{Our registered results: every success certified in both verdicts and long-tested at 128,
256 and 1{,}024 columns.}
\label{tab:ours}
\end{table}

``Nearly told'' is computed, not assumed. The sum $(A + B) \bmod p$ with the quotient output tells
apart 22 of division's 25 classes at $p = 5$ (44 of 49, 112 of 121 and 158 of 169 at 7, 11 and 13).
Only the $p - 1$ states with quotient $-1$, where $A \equiv -B \not\equiv 0$, share a label. Whether
T10's planted quotient and the 4-class output together determine a planted machine's state was not
computed. Its blocks are 4 of 16 states, so it is ``partly told'' at least.

\section{The claim, made strictly correct}

\begin{enumerate}
\item \textbf{Hidden state (stands, with a smaller number).} Among the work we found, the largest
  machines proved equal to their targets at every length, after gradient descent on outputs that
  leave the state hidden, are ours: the $S_4$ word problem and move-to-front on 4 items at 24 states
  (4 of 20 runs each), and random machines at 12 states (3 of 20) and, primitive, at 16 (2 of 30
  runs). The largest proved before is about 5 states
  \citep{zeng-1993-self-clustering-recurrent-networks,aichernig-2024-learning-minimal-automata-rnn}.
  The largest learned with the state hidden, though only tested, is 30
  \citep{umili-2024-deepdfa-automata-learning-probabilistic-relaxations}. The size is about five
  times the proved prior, at low rates: this is existence, not reliability.
\item \textbf{Told the state (no size claim).} With outputs that determine the state, our largest is
  division at 169 states (15 of 20). That is smaller than
  \citeauthor{giles-1995-learning-class-large-finite-state-machines}'s 512 and well below
  \citeauthor{brandoit-2026-length-independent-state-tracking}'s 7{,}920, so no size record is claimed
  here. What is ours is the controlled comparison: the same target and cell certify 16 of 20 told the
  state's coordinates and 0 of 20 on the quotient alone (121 states, Fisher $p = 7.7 \times
  10^{-8}$). The suffix control (state-carrying, no approximation, 0 of 20) says an approximation
  helps for being one, not merely for carrying state.
\item \textbf{Length.} Proof at every length is not ours alone:
  \citeauthor{brandoit-2026-length-independent-state-tracking} prove it too, by the same kind of
  enumeration. Ours applies it to Mealy machines whose outputs are not the state, with near ties
  branched both ways, and (K9) without the target in the loop.
\end{enumerate}

\section{Re-interpretation}

The correction moves the paper's weight from size to signal.
\begin{itemize}
\item \textbf{The 169 states are a result about approximations, not about size.} Division at 169
  states is learned when the network is told the state's coordinates. Told the quotient alone, it
  learns division at 25 states once in 20 and at 121 never. Paper~1's K4 and K6 read correctly as
  ``given its coordinates, a product machine of up to 169 states is learned reliably; without them,
  not at all''. That is a claim about what the signal must carry, and it is the one our data
  supports.
\item \textbf{The hidden-state frontier is about 24 states, and it is where the open problem is.} The
  algebras and random machines (T5, T10's P0) are the results nearest a ``larger than reported''
  claim. Their rates are 0--4 of 20, and random machines stop between 12 and 16 states. The
  machine-discovery aim (an unknown machine found from behaviour) lives here, where nothing tells the
  network the state.
\item \textbf{Approximations between the two.} The middle rows of Table~\ref{tab:ours} (the sum
  alone, nearly telling; T10's planted quotient, partly telling) are where approximations matter
  without handing over the state. T10b, P2 against P0 (8 against 2 machines, $p = 0.0206$, no
  margin), is the one registered test of an approximation that leaves the state partly hidden; T10r
  replicates it. T11 gives pairs that determine the state, so its rates, whatever they are, belong in
  the ``told'' class.
\item \textbf{What still has no precedent we found} is the reading of a trained network's
  decomposition against the target's congruences (K8), and the proof without the target in the loop
  (K9). Neither depends on a size record.
\end{itemize}

\section{Consequences for paper 1 (the first author's call)}

\begin{enumerate}
\item The paper workplan's ``result nobody else has'' (100--169 states, learned and proved) does not
  stand as worded. It should be restated as the controlled effect of the approximations (K4, K6), plus
  the hidden-state results at 24 states.
\item The abstract's ``up to 169 states'' should say the network was given the state's coordinates.
\item Related work should give \citeauthor{brandoit-2026-length-independent-state-tracking}'s sizes
  and supervision, and add \citet{giles-1995-learning-class-large-finite-state-machines},
  \citet{aichernig-2024-learning-minimal-automata-rnn} and DeepDFA's 30 tested states. The sentence
  ``the extracted-machine work stops at a few states'' is true only of hidden-state work with exact
  checks.
\item The introduction's question (``how large?'') should be asked per signal: how large when told
  the state, how large when not.
\end{enumerate}

\section{Limits of this check}

The search covered the repository's 300 papers and two web searches on 27 October 2026. A hidden-state
result above 24 states, proved, could exist outside them. Before the manuscript states the hidden-state
claim, it needs one more search aimed at exact extraction from networks trained on labels alone
(SAT-based or bisimulation checks, 2023--2026). The PDF of
\citeauthor{aichernig-2024-learning-minimal-automata-rnn} is not in the literature repository
(Springer refuses scripted downloads); its numbers here come from the publisher's full-text page.

\bibliographystyle{plainnat}
\bibliography{library}

\end{document}
